\documentclass[11pt,a4paper]{article}
\usepackage[margin=0.7in, top=0.6in, bottom=0.6in]{geometry}
\usepackage[utf8]{inputenc}
\usepackage[T1]{fontenc}
\usepackage{charter}
\usepackage{fontawesome5}
\usepackage{xcolor}
\usepackage{enumitem}
\usepackage{hyperref}
\usepackage{microtype}

\definecolor{accent}{HTML}{1E3A5F}     % Professional Dark Navy
\definecolor{accentblue}{HTML}{2B6CB0} % Slate Blue
\definecolor{lightgray}{HTML}{F3F4F6}
\definecolor{darktext}{HTML}{1F2937}
\definecolor{bulletline}{HTML}{9CA3AF}
\definecolor{subtletext}{HTML}{4B5563}

\hypersetup{colorlinks=true, linkcolor=accentblue, urlcolor=accentblue}

\setlength{\parindent}{0pt}
\setlength{\parskip}{8pt}

\begin{document}
\thispagestyle{empty}
\color{darktext}

% Encabezado idéntico al CV
\begin{center}
    {\Huge \bfseries\color{accent} SANTIAGO ALEJANDRO OCHOA QUESADA} \\[6pt]
    {\small\color{darktext}
    \faEnvelope\ \href{mailto:santy.alejo4@gmail.com}{santy.alejo4@gmail.com}
    \quad\textcolor{bulletline}{\textbar}\quad
    \faPhone\ +57 320 351 7153
    \quad\textcolor{bulletline}{\textbar}\quad
    \faMapMarker*\ Bogotá, Colombia \\[4pt]
    \faLinkedin\ \href{https://linkedin.com/in/santiago-alejandro-ochoa-quesada-a07a18288/}{LinkedIn}
    \quad\textcolor{bulletline}{\textbar}\quad
    \faGithub\ \href{https://github.com/Sochoaq}{GitHub}
    }
\end{center}
\vspace{2pt}
{\color{bulletline}\hrule height 0.8pt}
\vspace{10pt}

% Metadatos de la carta
\noindent
\begin{minipage}[t]{0.62\textwidth}
    \textbf{\color{accent} Para:}\\
    Equipo de Selección \& Talento Humano\\
    \textbf{Quala S.A.}\\
    Bogotá, Colombia
\end{minipage}
\hfill
\begin{minipage}[t]{0.35\textwidth}\raggedleft
    \textbf{\color{accent} Fecha:}\\
    22 de octubre de 2026
\end{minipage}

\vspace{6pt}
\noindent
{\large\bfseries\color{accent} Asunto: Postulación al cargo de Ingeniero(a) de Proyectos de Seguridad Tecnológica - TI}
\vspace{2pt}

\noindent Estimado Equipo de Selección \& Talento Humano,

Presento mi postulación para el rol de \textbf{Ingeniero(a) de Proyectos de Seguridad Tecnológica - TI} en Quala. Como Ingeniero Electrónico de la \textbf{Universidad Nacional de Colombia}, con un promedio de 4.6/5.0 y sólida experiencia actual como \textbf{Ingeniero de Proyectos} en AM4, poseo la combinación exacta de rigor técnico, gestión de infraestructura y capacidad de innovación para fortalecer la postura de ciberseguridad y continuidad del negocio en su organización.

A lo largo de mi trayectoria, he liderado de principio a fin el ciclo de vida de proyectos tecnológicos complejos, abarcando desde la arquitectura de sistemas bajo Linux y redes de comunicación hasta la implementación de automatizaciones robustas en Python y C/C++. Mi enfoque no solo se limita al cumplimiento estricto de cronogramas y estándares de calidad, sino también a la incorporación estratégica de Inteligencia Artificial y análisis de datos para la detección temprana de riesgos y la protección proactiva de la información.

Residiendo actualmente en Bogotá y con disponibilidad para el modelo híbrido, me entusiasma profundamente la oportunidad de aportar mi experiencia en gestión técnica, infraestructura y automatización en Quala, impulsando la evolución de la seguridad tecnológica y garantizando soluciones resilientes y de alto impacto.

Agradezco de antemano su atención y consideración. Quedo a su completa disposición para profundizar en cómo mis competencias técnicas y de liderazgo pueden sumar valor inmediato a su equipo.

\vspace{10pt}
\noindent Cordialmente,

\vspace{12pt}
\noindent
\textbf{\large\color{accent} Santiago Alejandro Ochoa Quesada}\\
{\small\color{subtletext} Ingeniero Electrónico | AI Engineering \& Software}\\
{\footnotesize\color{subtletext} santy.alejo4@gmail.com \;\textbar\; +57 320 351 7153 \;\textbar\; Bogotá, Colombia}

\end{document}
